
This study asked whether static analyzers can improve LLM code repair. The answer remains open---but the methodology to find it is now established.

\textbf{Findings:} Canonical HumanEval solutions exhibit only 9.15\% pylint warning rates (15 of 164 problems). This makes them unsuitable proxies for LLM-generated code. The MUST\_WORK gate correctly identified this limitation before downstream experiments consumed resources.

\textbf{Artifacts:} A validated pylint analysis pipeline for code generation research, reusable components for HumanEval loading and static analysis, and a gate-based hypothesis validation framework.

\textbf{Methodology insight:} Gate-based validation functions as intended. By testing existence assumptions before full experiments, methodology issues are caught early.

\textbf{Next steps:} Re-run the existence gate with actual LLM generations. If warning rates exceed 30\%, proceed to mechanism validation and full pass@k comparison.

\textbf{Broader implication:} Before evaluating any LLM intervention, validate that the intervention has signal. Canonical solutions are not LLM output---a lesson applicable beyond static analysis to any benchmark-based evaluation.

The question of whether static analysis improves pass@k merits investigation. This work provides the validated tools to pursue that investigation.
